% Lemma 3.2.2 — the localization at the unique prime over p
% Blueprint node for proofs/lem-3-2-2.md. Statement carries the standing
% hypothesis p | f explicitly (ERRATA E6.2); it is FALSE for split p.

\begin{lemma}[Localization at the unique prime over $p$]
  \label{lem-3-2-2}
  \lean{QuadraticOrder.exists_int_denominator,
    QuadraticOrder.isUnit_algebraMap_atPrime_iff_not_dvd_norm,
    QuadraticOrder.existsUnique_isPrime_mem_of_dvd_conductor}
  \uses{notation, prop-3-2-1}
  Let $d = Df^2 < 0$ be an imaginary quadratic discriminant with $D$ fundamental
  and $f \ge 1$ the conductor, and let
  $\mathcal{O} = \mathcal{O}_d = \mathbb{Z}[\tau]$, $\tau = (d + \sqrt{d})/2$,
  with $g(x) = x^2 - dx + \frac{d^2 - d}{4}$ the minimal polynomial of $\tau$,
  $N(\alpha) = \alpha\bar\alpha$ the norm, and
  $\mathbb{Z}_{(p)} = \{a/t : a, t \in \mathbb{Z},\ p \nmid t\}$.
  Assume $p$ is a rational prime with $p \mid f$. Then:
  \begin{enumerate}
    \item[(0)] $\mathcal{O}$ has exactly one prime ideal $\mathfrak{P}$ lying
      over $p$, namely $\mathfrak{P} = (p, \tau - A_0)$ where $x - A_0$ is the
      unique monic irreducible factor of $g \bmod p$ ($A_0 = 0$ for odd $p$;
      $A_0 \equiv d/4 \pmod 2$ for $p = 2$); in particular $p$ is not split in
      $\mathcal{O}$ and $\mathcal{O}/\mathfrak{P} \cong \mathbb{F}_p$.
    \item[(i)] (\emph{Integer denominators.}) Every $z$ in the localization
      $\mathcal{O}_{\mathfrak{P}} = \{\alpha/\beta : \alpha \in \mathcal{O},\
      \beta \in \mathcal{O} \setminus \mathfrak{P}\} \subseteq K = \mathbb{Q}(\sqrt{d})$
      can be written $z = \alpha/t$ with $\alpha \in \mathcal{O}$,
      $t \in \mathbb{Z}$, $p \nmid t$; equivalently, inside $K$,
      \[
        \mathcal{O}_{\mathfrak{P}}
          = \mathbb{Z}_{(p)}[\tau]
          = \mathbb{Z}_{(p)} + \mathbb{Z}_{(p)}\tau .
      \]
    \item[(ii)] (\emph{Unit criterion.}) For $\alpha = x + y\tau \in \mathcal{O}$,
      the image of $\alpha$ in $\mathcal{O}_{\mathfrak{P}}$ is a unit if and
      only if $p \nmid N(\alpha) = x^2 + dxy + \frac{d^2 - d}{4}y^2$ in
      $\mathbb{Z}$.
  \end{enumerate}
\end{lemma}

\begin{proof}
  \uses{prop-3-2-1}
  \textbf{Step 0 (uniqueness of the prime over $p$).}
  By Proposition~\ref{prop-3-2-1}, the primes of $\mathcal{O}$ containing $p$
  are exactly the ideals $p\mathcal{O} + (g_i(\tau))$ where $\bar g_i$ runs over
  the monic irreducible factors of $\bar g := g \bmod p$ in $\mathbb{F}_p[x]$.
  We claim $\bar g$ is the square of a monic linear polynomial. If $p$ is odd,
  $p \mid f$ gives $p^2 \mid d$, and
  $v_p\bigl(\frac{d^2-d}{4}\bigr) = v_p(d(d-1)) = v_p(d) \ge 2$ since $v_p(4)=0$
  and $p \nmid d-1$; hence $\bar g = x^2$. If $p = 2$, then $2 \mid f$ gives
  $4 \mid f^2 \mid d$, the linear term $-dx$ vanishes mod $2$, and the constant
  term satisfies $\frac{d^2-d}{4} = \frac{d}{4}(d-1) \equiv \frac{d}{4} \pmod 2$
  because $d - 1$ is odd; hence $\bar g = x^2$ when $8 \mid d$ and
  $\bar g = x^2 + 1 = (x+1)^2$ when $v_2(d) = 2$. In all cases
  $\bar g = (x - \bar A_0)^2$ has a single monic irreducible factor, so
  $\mathfrak{P} = (p, \tau - A_0)$ is the unique prime of $\mathcal{O}$
  containing $p$, with $\mathcal{O}/\mathfrak{P} \cong
  \mathbb{F}_p[x]/(x - \bar A_0) \cong \mathbb{F}_p$. Any prime containing $p$
  lies over $p$: $\mathfrak{P} \cap \mathbb{Z}$ is a proper ideal of
  $\mathbb{Z}$ containing the maximal ideal $p\mathbb{Z}$, hence equals
  $p\mathbb{Z}$. This proves (0); in particular $p$ is not split in
  $\mathcal{O}$.

  \textbf{Step 1 ($\mathfrak{P}$ is stable under conjugation).}
  Conjugation $\iota : \alpha \mapsto \bar\alpha$ restricts to a ring
  automorphism of $\mathcal{O}$, since $\iota(x + y\tau) = (x + yd) - y\tau
  \in \mathcal{O}$ and $\iota$ is an involutive field automorphism of $K$.
  Thus $\iota(\mathfrak{P})$ is a prime ideal of $\mathcal{O}$ containing
  $\iota(p) = p$, and the uniqueness of Step~0 gives
  $\iota(\mathfrak{P}) = \mathfrak{P}$.

  \textbf{Step 2 (norm detects membership in $\mathfrak{P}$).}
  We claim: for $\alpha \in \mathcal{O}$,
  $\alpha \in \mathfrak{P} \iff p \mid N(\alpha)$.
  ($\Rightarrow$) If $\alpha \in \mathfrak{P}$ then
  $\bar\alpha \in \iota(\mathfrak{P}) = \mathfrak{P}$ by Step~1, so
  $N(\alpha) = \alpha\bar\alpha \in \mathfrak{P} \cap \mathbb{Z} = p\mathbb{Z}$.
  ($\Leftarrow$) If $p \mid N(\alpha)$ then
  $\alpha\bar\alpha \in p\mathbb{Z} \subseteq \mathfrak{P}$; since
  $\mathfrak{P}$ is prime, $\alpha \in \mathfrak{P}$ or
  $\bar\alpha \in \mathfrak{P}$, and in the latter case applying $\iota$ gives
  $\alpha = \iota(\bar\alpha) \in \iota(\mathfrak{P}) = \mathfrak{P}$.

  \textbf{Step 3 (proof of (ii)).}
  Since $\mathcal{O}_{\mathfrak{P}} \subseteq K$ with $K$ a field, an element of
  $\mathcal{O}_{\mathfrak{P}}$ is a unit iff it is nonzero and its inverse in
  $K$ lies in $\mathcal{O}_{\mathfrak{P}}$. If $\alpha \notin \mathfrak{P}$,
  then $\alpha \ne 0$ and $1/\alpha \in \mathcal{O}_{\mathfrak{P}}$ by
  definition, so $\alpha$ is a unit. If $\alpha \in \mathfrak{P}$ were a unit,
  write $1/\alpha = \gamma/\delta$ with $\gamma \in \mathcal{O}$,
  $\delta \notin \mathfrak{P}$; then $\delta = \alpha\gamma \in \mathfrak{P}$,
  a contradiction. Hence $\alpha$ is a unit in $\mathcal{O}_{\mathfrak{P}}$
  iff $\alpha \notin \mathfrak{P}$, which by Step~2 holds iff
  $p \nmid N(\alpha)$. Expanding $N(x + y\tau)$ via
  $\tau + \bar\tau = d$ and $\tau\bar\tau = \frac{d^2-d}{4}$ gives the stated
  form.

  \textbf{Step 4 (proof of (i)).}
  If $t \in \mathbb{Z}$ with $p \nmid t$, then
  $t \notin \mathfrak{P}$ (as $\mathfrak{P} \cap \mathbb{Z} = p\mathbb{Z}$),
  so $\alpha/t \in \mathcal{O}_{\mathfrak{P}}$ for all $\alpha \in \mathcal{O}$;
  this inclusion holds for any prime over any $p$. Conversely, let
  $z = \alpha/\beta \in \mathcal{O}_{\mathfrak{P}}$ with $\alpha \in \mathcal{O}$
  and $\beta \in \mathcal{O} \setminus \mathfrak{P}$. By Step~2,
  $p \nmid N(\beta)$, and $N(\beta) \neq 0$ since $\beta \neq 0$. Then
  \[
    z = \frac{\alpha}{\beta}
      = \frac{\alpha\bar\beta}{\beta\bar\beta}
      = \frac{\alpha\bar\beta}{N(\beta)},
  \]
  with numerator $\alpha\bar\beta \in \mathcal{O}$ and integer denominator
  $t = N(\beta)$ prime to $p$. Finally, bringing coordinates to a common
  denominator,
  $\{\alpha/t : \alpha \in \mathcal{O},\, t \in \mathbb{Z},\, p \nmid t\}
    = \mathbb{Z}_{(p)} + \mathbb{Z}_{(p)}\tau$,
  and $\mathbb{Z}_{(p)}[\tau] = \mathbb{Z}_{(p)} + \mathbb{Z}_{(p)}\tau$
  because $\tau$ satisfies the monic quadratic $g$ over
  $\mathbb{Z} \subseteq \mathbb{Z}_{(p)}$.
\end{proof}

\paragraph{Remark (failure for split $p$; ERRATA E6.2).}
Both (i) and (ii) are \emph{false} without the hypothesis $p \mid f$ (more
precisely, without the uniqueness of the prime over $p$ that it forces). For
$d = -4$, $\mathcal{O} = \mathbb{Z}[i]$, $p = 5 = (2+i)(2-i)$,
$\mathfrak{P} = (2+i)$: the element $2 - i \notin \mathfrak{P}$ is a unit of
$\mathcal{O}_{\mathfrak{P}}$ although $5 \mid N(2-i) = 5$, violating (ii); and
$1/(2-i) = (2+i)/5 \in \mathcal{O}_{\mathfrak{P}}$ admits no integer
denominator prime to $5$ (from $t(2+i) = 5\alpha$ one gets $t = (2-i)\alpha$,
whence $t^2 = 5N(\alpha)$ and $5 \mid t$), violating (i). The thesis states the
lemma for an arbitrary prime with the hypothesis ``$p$ not split in
$\mathcal{O}$''; following ERRATA E6.2 the blueprint version carries the
standing hypothesis $p \mid f$ under which non-splitness is proved (Step~0)
rather than assumed. The thesis's original proof of Step~2 factors
$N(\beta) = [\mathcal{O} : \beta\mathcal{O}]$ into local indices over the
primary components of $\beta\mathcal{O}$; the conjugation argument above is an
equivalent, formalization-friendly replacement resting on the same pivot
(uniqueness of the prime over $p$).
